% Folder 57, batch 04 (archivists' pages 61-80).
% Pass: Opus 5.5 (claude-opus-5-5), 2026-10-03 — first pass, unchecked against the pages by a human.
%
% Pages 64, 66 and 79 are covers inscribed in his hand; their words become
% section headings and the pages get no \page{}. Pages 72-78 are his own letter
% to Murre of 18 July 1962 (two typed states: a fair copy, pp. 72-74, and the
% corrected draft, pp. 75-78), transcribed as his own correspondence.
\documentclass[11pt,a4paper]{article}
\input{../preamble/grothendieck}

\folder{57}
\batch{4}
\pages{61}{80}
\foldertitle{Picard : tapuscrit annoté (s.d.), lettres (s.d., 1962).}
\dating{1962-[vers 1968]}
\watermark{Édition de démonstration}

\begin{document}

\section{Schéma d'Albanese et schéma de Picard}

\page{61}
\note{La page commence au milieu d'un argument entamé avant le lot : il s'agit de la dualité entre $\underline{\mathrm{Pic}}_{X/S}$ et le schéma d'Albanese. Main très rapide ; la prose de liaison n'est souvent lisible qu'en partie.}
\ill{} vérifie aisément \uncertain{que pour} \uncertain{tout} schéma abélien $B/S$ \ill{} a une bijection
\[
\mathrm{Hom}_{S\text{-gr}}(\hat{B}, \underline{\mathrm{Pic}}_{X/S}) \xleftarrow{\ \sim\ } \mathrm{Hom}_{S\text{-gr}}(B, \underline{B}_{X/S})
\]
\note{le premier argument, $\hat{B}$, est écrit sur un $A$ biffé}
d'où par \emph{dualité} \emph{des} V.A. une bijection
\[
\begin{cases}
\mathrm{Hom}_{S\text{-gr}}(\hat{A}, \underline{\mathrm{Pic}}_{X/S}) \xleftarrow{\ \sim\ } \mathrm{Hom}_{S\text{-gr}}(\underline{A}^{0}_{X/S}, A) \\[4pt]
\text{où } \underline{A}^{0}_{X/S} = \hat{\underline{B}}_{X/S}
\end{cases}
\]

Donc pour \emph{un schéma abélien $A$ sur $S$}, \emph{\uncertain{la donnée}} \emph{d'un fibré principal hom. $P$ \uncertain{sous} $A$ et d'une $S$-\uncertain{morphisme} $\psi\colon X \to P$} \emph{\uncertain{équivaut} \ill{}} \emph{à} \add{à $S$-gr} \emph{la donnée d'un hom.} $u\colon \underline{A}^{0}_{X/S} \to A$. \ill{} \ill{} \uncertain{précis}, il y a un fibré principal hom. \add{et un $S$-\uncertain{morphisme}} \struck{\ill{}} $\underline{A}^{1}_{X/S}$ sous $\underline{A}^{0}_{X/S}$, canonique, $\psi_{X/S}\colon X \to \underline{A}^{1}_{X/S}$, et la \uncertain{correspondance} \uncertain{explicite} \uncertain{dessus} \ill{}, s'obtient \uncertain{en} \uncertain{prenant} :
\[
\text{\struck{$X$}}\quad P \simeq \underline{A}^{1}_{X/S} \times^{\underline{A}^{0}_{X/S}} (A, u)\ \ldots\ldots
\]
\marginal{// \uncertain{On} \uncertain{notera}} Les schémas d'Albanese $\underline{A}^{0}_{X/S}$ et $\underline{A}^{1}_{X/S}$ \uncertain{resp.} \emph{homogènes} et \emph{\uncertain{inhomogènes}} \uncertain{de} $X/S$. Ils \uncertain{sont} \uncertain{définis} \uncertain{à} \uncertain{unique} \uncertain{isom.} \uncertain{près}. \emph{\uncertain{Le} \uncertain{système} \uncertain{triple}} $(\underline{A}^{0}_{X/S}, \underline{A}^{1}_{X/S}, \psi_{X/S})$ \emph{représente le foncteur} $T_X(P)$ \uncertain{par} \uncertain{construction}.

\page{62}
\uncertain{Les} \uncertain{notions} \ill{} \ill{} \uncertain{temps} \uncertain{que} \uncertain{par} \emph{construction}, \struck{\ill{}} la formation de $\underline{A}^{0}_{X/S}$, $\underline{A}^{1}_{X/S}$, $\psi_{X/S}$ \uncertain{commute} \uncertain{au} \uncertain{changement} de base \uncertain{quelconque}.

[N.B. Il y a \uncertain{peut-être} \uncertain{lieu} de \uncertain{définir} \uncertain{axiomatiquement} un \emph{schéma d'Albanese} de $X/S$ par la condition de \uncertain{représenter} le foncteur $T_X(P)$, $P$ \uncertain{parcourant} \uncertain{les} schémas abéliens, \uncertain{mais} \uncertain{d'une} $X$-\ill{} et \ill{} \uncertain{fait} que les \ill{} \uncertain{deux} \uncertain{soient} \uncertain{stables} \uncertain{par} le changement de base. Pour \emph{l'existence}, il \emph{suffit} \ill{} \ill{} \uncertain{nécessairement} \ill{}, \uncertain{lorsque} \uncertain{par} \uncertain{l'une} \uncertain{précisément} \ill{} \uncertain{un} \uncertain{sous-schéma} abélien de $\underline{\mathrm{Pic}}_{X/S}$ \uncertain{ayant} $\underline{\mathrm{Pic}}^{0}_{X/S}$ \ill{} \ill{} \ill{}. \uncertain{Ainsi}, il \ill{} \uncertain{plausible} \uncertain{qu'un} schéma d'Albanese \uncertain{existe} \uncertain{pour} $X/S$ \uncertain{quand} \ill{} $S$ \uncertain{est} le \uncertain{spectre} d'un corps $k$, \ill{} \ill{} $\underline{\mathrm{Pic}}^{0}_{X/k}$ n'est pas propre sur $k$ \add{\ill{}} [\uncertain{plus} \uncertain{grand} \uncertain{sous-schéma} abélien de $\underline{\mathrm{Pic}}_{X/k}$ ; il \uncertain{faut} \uncertain{prouver} \uncertain{évidemment} \uncertain{une} propriété universelle de \uncertain{ce} \uncertain{dernier}, \uncertain{par} \uncertain{les} changements de base …]].

\[
\begin{cases}
X \longrightarrow \underline{A}^{1}_{X/S} \\[4pt]
\underline{\mathrm{Pic}}^{0}_{X/S} \longleftarrow \underline{\mathrm{Pic}}^{0}_{\underline{A}^{1}/S} \simeq \underline{\mathrm{Pic}}^{0}_{\underline{A}^{0}/S}
\end{cases}
\]
\drawing{sous ce système, une flèche montante $\underline{B} \to \underline{\mathrm{Pic}}^{0}_{X/S}$ et une flèche oblique, barrée d'un trait, de $\underline{\mathrm{Pic}}^{0}_{\underline{A}^{1}/S}$ vers $\underline{B}$ ; le croquis s'arrête au bas de la page}

\page{63}
N.B. \struck{\uncertain{Question}} \emph{\uncertain{Question}}\note{le mot souligné est surchargé : peut-être « Question » réécrit} Si $\underline{\mathrm{Pic}}_{X/S}$ est simple sur $S$ \uncertain{avec} \uncertain{fibres} \uncertain{des} $\underline{\mathrm{Pic}}^{0}_{X/S}$ [\ill{} \struck{\ill{}} \ill{} \ill{} \ill{}, \ill{} \uncertain{le} \uncertain{long} \uncertain{de} \uncertain{la} \uncertain{section} \uncertain{unité} …] \struck{\uncertain{les} \uncertain{fibres} \uncertain{des} \ill{} \ill{} \ill{} $S$, \ill{} \uncertain{propres}} \struck{Alors} alors $\underline{\mathrm{Pic}}^{0}_{X/S}$ est \emph{ouvert} dans $\underline{\mathrm{Pic}}_{X/S}$ [\uncertain{c'est} \uncertain{bien} \uncertain{connu}. \uncertain{En} \uncertain{effet}, \ill{} loc. de type fini sur $S$, car s'il est \emph{\uncertain{simple}} sur $S$, alors \uncertain{la} \uncertain{composante} \uncertain{neutre} \uncertain{des} \uncertain{fibres} \uncertain{forment} \uncertain{par} \uncertain{les} \uncertain{composantes} \ill{} \uncertain{une} \uncertain{partie} \uncertain{des} \ill{} \marginal{IV 7.(6.6 \struck{\ill{}})} \uncertain{fibres} \uncertain{géom.} \emph{loc. intègres} \struck{\ill{}} $G^{0}$ est \uncertain{ouvert} \uncertain{dans} $G$].
Supposons \uncertain{le} \uncertain{propre} sur $S$, \uncertain{donc} \uncertain{c'est} \uncertain{un} \uncertain{schéma} \uncertain{abélien}.
\uncertain{Alors} (a) les fibres de $\underline{\mathrm{Pic}}_{X/S}$ sont simples ;
(b) $\mathcal{O}_X$ est \uncertain{coh.} plat sur $S$ en \uncertain{codimension} \emph{1}, i.e. $\underline{\omega}_{\underline{\mathrm{Pic}}_{X/S}}$ est loc. libre …

[\struck{\uncertain{Réc.}} \uncertain{C'est-il} vrai \uncertain{réciproquement} que (a) et (b) \uncertain{impliquent} que $\underline{\mathrm{Pic}}_{X/S}$ est simple sur $S$ le long de la section unité ??? Notons qu'on \uncertain{a} \uncertain{évidemment} \ill{} (a) est \uncertain{automatiquement} vrai, et \struck{\ill{}} je pense que (b), dans le \uncertain{cadre} plus général \uncertain{où} $G$ est un \uncertain{groupe} \uncertain{sur} $S$ \struck{\ill{}}, $S$ \uncertain{réduit}, \uncertain{à} \uncertain{fibres} \uncertain{lisses}, et si $\underline{\omega}_{G/S}$ [\uncertain{formé} \uncertain{le} \uncertain{long} \uncertain{de} la section \uncertain{unité}] est \emph{loc. libre}, \emph{alors} $G$ \uncertain{est} \uncertain{simple} le long de la section unité …

De plus, il y a \uncertain{lieu} de \uncertain{se} \uncertain{demander} si \uncertain{pour} $X/S$ \uncertain{propre} et \uncertain{simple}, et \uncertain{de} \uncertain{car.} $0$, \uncertain{la} \uncertain{condition} \uncertain{générique} \uncertain{en} \uncertain{dim} \uncertain{est} \uncertain{satisfaite} \uncertain{automatiquement}. \uncertain{On} \uncertain{a} \uncertain{une} \uncertain{démonstration} \uncertain{purement} \emph{\uncertain{algébrique}} \uncertain{via} $\underline{\mathrm{Pic}}$ dans le cas où $S$ est \emph{réduit}. \uncertain{Dans} le cas \uncertain{où} $S$ \uncertain{est} \uncertain{artinien}, je \uncertain{ne} \uncertain{sais} \uncertain{rien}, \emph{\uncertain{au} \uncertain{contraire}} !]

\section{Plan}
\note{p.~64 : feuillet de couverture, de sa main, portant le seul mot « Plan ».}

\page{65}
\emph{Picard}

\begin{enumerate}
\item[1.] Théorème d'existence absolus
\item[2.] Théorème d'existence relatifs.
\item[3.] Théorème \add{Critères d'équivalence et} de finitude. [\uncertain{Néron}-\uncertain{Igusa} exclus ??]
\item[4.] $\underline{\mathrm{Pic}}^{\tau}_{X/S}$, équivalence numérique, gr.\ de Néron-Severi.
\item[\struck{6}5.] Critères de \struck{\ill{}} \uncertain{séparation}, de propreté pour $\underline{\mathrm{Pic}}_{X/S}$.
\item[\struck{6}.] Critères de \uncertain{simplicité} pour $\underline{\mathrm{Pic}}_{X/S}$
\item[\struck{7}8.] Albanese
\item[\struck{7}.] Dualité schémas abéliens. [\uncertain{notamment} \uncertain{auto-dualité}.]
\item[9.] \emph{Picard} des variétés abéliennes [\uncertain{Torelli} ?]
\item[\struck{6}5.] \struck{\ill{}} Variations sur la cour \add{\uncertain{correspondances} divisorielles} et la cube [\uncertain{multiplication} \uncertain{scalaire} $\underline{\mathrm{Pic}}^{\tau}_{X/S}$]
\item[10.] Compléments au dessus d'un corps.
[\uncertain{notamment} \uncertain{cas} des variétés \uncertain{non} complètes, des diviseurs \uncertain{sur} \uncertain{les} \uncertain{courbes}, \uncertain{le} \uncertain{groupe} \uncertain{de} une birationnelle, \struck{\ill{}} critères d'équivalence \uncertain{des} \ill{}, le \uncertain{cas} \uncertain{normal} \uncertain{etc} \struck{\ill{}} ; \uncertain{aussi} \uncertain{rappel} \uncertain{des} \uncertain{résultats} spéciaux obtenus \uncertain{dans} \uncertain{les} \uncertain{p.} \uncertain{précédentes} \uncertain{antérieures}]
\item[11.] Phénomènes spéciaux en car.\ 0.
\item[12.] (?) Picard locaux et progressions …
\end{enumerate}
\marginal{\uncertain{utiliser} \uncertain{les} \uncertain{variétés} \uncertain{des} \uncertain{schémas} \uncertain{en} \uncertain{groupes} …}
\note{La numérotation est remaniée à l'encre : des flèches en marge reportent l'ancien 6 (critères de simplicité) entre 4 et 5, l'ancien 6 (variations sur le carré et le cube) devient 5, Albanese passe de 7 à 8, la dualité de 8 à 7. Les numéros biffés sont donnés tels qu'on les lit ; l'ordre définitif n'est pas sûr.}

Au bas de la page, un schéma des dépendances entre chapitres :

\begin{tikzcd}[column sep=small, row sep=small]
 & 1 \arrow[dl] \arrow[dr, Rightarrow] & 2 \arrow[d, Rightarrow] \\
\cdot \arrow[dd, dashed] & & 3 \arrow[d, Rightarrow] \\
 & & 4 \arrow[dll, dashed] \arrow[d, Rightarrow] \\
6 \arrow[drr, Rightarrow] & & 5 \arrow[d, Rightarrow] \\
 & & 7 \arrow[d, Rightarrow] \\
 & & 8 \arrow[d, Rightarrow] \\
 & & 9 \arrow[r] & 10
\end{tikzcd}

\note{Redessiné d'après le croquis : de 1 partent une flèche simple vers la gauche et une double vers 3 ; une flèche pointillée descend de la gauche vers 6, marquée d'un « 5 » \uncertain{biffé} ; une pointillée va de 4 à 6 ; une double flèche oblique va de 6 vers 7. Après le 9, un mot biffé illisible précède la flèche vers 10. La colonne $2 \Rightarrow 3 \Rightarrow 4 \Rightarrow 5 \Rightarrow 7 \Rightarrow 8 \Rightarrow 9 \to 10$ est sûre ; la place exacte des flèches obliques l'est moins.}

\section{Théorème de finitude}
\note{p.~66 : feuillet de couverture, de sa main, « Th. de finitude », souligné.}

\page{67}
\marginal{\uncertain{Finitude} de $\underline{\mathrm{Pic}}^{\tau}_{X/S}$ — au crayon} \marginal{(\emph{Dém} de Matsusaka) — au crayon, encadré}
Soit $X$ une surface projective non singulière, sur un corps algébriquement clos $k$. \add{Soit $H$ un diviseur ample.} Rappelons la formule de R.-R. faible
\[
\chi(D) = \tfrac12 D\cdot(D-K) + c
\]
où $c$ est une constante. \struck{Soit $D'$ numériquement \add{Comme} équivalent à $0$, on trouve}
\[
\chi(D) = \ell(D) + \ell(K-D) - \dim H^1(X, \mathcal{L}(D))
\]
d'où
\[
\text{\struck{$\chi(\ell$}}\quad \text{\struck{$\chi(D)$}}\quad
\ell(D) + \ell(K-D) \geqslant \tfrac12 D\cdot(D-K) + c
\]
\struck{Donc on en conclut}
Supposons p.~ex.\ que $\ell(K-D) = 0$, ce qui est le cas si $H\cdot(K-D) \leqslant 0$, i.e.\ si
\[
\text{a)}\quad \text{\struck{\ill{}}}\ (D-K)\cdot H \geqslant 0
\]
Alors on a $\ell(D) \geqslant \frac12 D(D-K) + c$, donc si \add{et} \struck{\ill{}}
\[
\text{b)}\quad \tfrac12 D(D-K) + c \geqslant 0
\]
\uncertain{implique} que $\ell(D) > 0$. Les conditions a) et b) sont satisfaites p.~ex.\ pour \struck{\ill{}} $D = nH$, $n$ assez grand. Si alors $\Delta$ \struck{\ill{}} est un diviseur numériquement équivalent à $0$, les conditions a) et b) \uncertain{sont} vérifiées par $D$, donc on a
\[
\ell(D + \Delta) > 0
\]
donc $D + \Delta$ est représenté par un diviseur positif. Donc :

\page{68}
\emph{Proposition} Soit $D$ un diviseur satisfaisant a) et b), et soit $\mathfrak{M}$ la famille des \emph{diviseurs} $\geqslant 0$ \emph{numériquement} équivalents à $D$. Alors \struck{\ill{}} \uncertain{tout} diviseur numériquement équivalent à $0$ est lin.\ équivalent à un diviseur de la forme $D' - D$, avec $D' \in \mathfrak{M}$, \uncertain{donc} \uncertain{i.e.} $\mathfrak{M}$ \uncertain{est} une classe \uncertain{dans} $\underline{\mathrm{Pic}}^{\mathfrak{n}}_{X/k}$.

Le th.\ de \struck{finitude} R.R. \add{(faible) sur $X$} \uncertain{montre} que le polynôme de Hilbert de $X \in \mathfrak{M}$ \uncertain{est} \uncertain{borné} \uncertain{ou} \uncertain{fixé}. Le th.\ de finitude de Chow (\uncertain{appliqué} \uncertain{aux} \uncertain{schémas} de Hilbert) montre donc que $\mathfrak{M}$ est quasi-compact.
\marginal{\uncertain{plus} \uncertain{simplement}, les degrés des diviseurs \uncertain{de} $\mathfrak{M}$ \uncertain{restant} fixés …}
Donc

\emph{Corollaire} $\underline{\mathrm{Pic}}^{n}_{X/k}$ est qu.\ compact \struck{\ill{}}

\uncertain{Cela} équivaut à :

\emph{Corollaire} $\underline{\mathrm{Pic}}^{n}_{X/k} / \underline{\mathrm{Pic}}^{0}_{X/k}$ est fini

ou encore

\emph{Corollaire} $\underline{\mathrm{Pic}}^{n}_{X/k} = \underline{\mathrm{Pic}}^{\tau}_{X/k}$, le groupe de torsion \struck{\ill{}} dans $\mathrm{NérSév}_{X/k}$ est fini.

La relation $\underline{\mathrm{Pic}}^{\tau}_{X/k} = \underline{\mathrm{Pic}}^{\tau}_{X/k}$ \uncertain{signifie} aussi :\note{la première lettre en exposant est peut-être un trait et non un $\tau$}

\emph{Corollaire} Soit $N = \underline{\mathrm{Pic}}_{X/k} / \underline{\mathrm{Pic}}^{\tau}_{X/k}$. Alors la forme d'intersection sur $N$ est non dégénérée.

\page{69}
Notons dans ce contexte le résultat plus profond de \uncertain{Igusa} :

\emph{Proposition}. $X$ comme ci-dessus, \struck{\ill{}} posons
\[
E(X) = \sum_{i=0}^{4} (-1)^i\, b_i(X)
\]
avec
\[
\begin{cases}
b_0(X) = b_4(X) = 1 \\
b_1(X) = b_3(X) = \dim \underline{\mathrm{Pic}}_X \\
E(X) = c_2(X) = \text{self-intersection de la diagonale dans } X\times X.
\end{cases}
\]
\uncertain{On} \uncertain{peut} \uncertain{déterminer} donc $b_2(X)$ de façon \uncertain{numérique}, \add{$b_2 = E - 2 + 2p$ où $p = \dim \underline{\mathrm{Pic}}_X$} \struck{\uncertain{c'est-à-dire}, \ill{}}. \uncertain{Si}
\[
\rho = \mathrm{rg}\,(\mathrm{NérSév}_{X/k}),
\]
\uncertain{alors} \uncertain{on} \uncertain{a} l'inégalité
\[
\boxed{\rho \leqslant b_2}
\]

\section{Théorème de finitude sur Picard (Commentaire)}

\page{70}
\emph{Théorème de finitude sur Picard} \marginal{(Commentaires — au crayon}

Kleiman prouve, pour $X/k$ \emph{projectif}, les critères (viii)\note{« viii » cerclé ; de même plus bas (v), (vi), (vii)} de $\tau$-équivalence à $0$, sans utiliser le fourbis délicat (i) (ii) …. Pour récupérer le cas $X/k$ propre, il doit utiliser (i) (via lemme de Chow). Avec le même grain de sel, il prouve (v) que $\underline{\mathrm{Pic}}^{\tau}_{X/S}$ est qu.-cpt.\ sur $S$. [\emph{En fait, il prouve directement que} $\underline{\mathrm{Pic}}^{\mathrm{num}}_{X/S}$ \emph{est quasi-compact} …].

\marginal{||} Utilisant le th.\ de l'index pour une surface projective non singulière, Kleiman prouve de plus le « critère d'équivalence » sous la forme suivante : si $X$ est projectif$/k$, \uncertain{toutes} ses composantes de dim $\geqslant 3$, et $Y$ section hyperplane, d'où $f^{*}\colon \underline{\mathrm{Pic}}_{X/k} \to \underline{\mathrm{Pic}}_{Y/k}$, \note{l'indice de $\underline{\mathrm{Pic}}_{X/k}$ est écrit sur un $Y$ biffé} alors $f^{*-1}(\underline{\mathrm{Pic}}^{\tau}_{Y/k}) = \underline{\mathrm{Pic}}^{\tau}_{X/k}$, i.e.\ \uncertain{si} \uncertain{le} \uncertain{hom.} \uncertain{des}
$\mathrm{NSL}(X) \to \mathrm{NSL}(Y)$ est injectif. \struck{Ceci} \add{Ceci, qui précède}, et l'\uncertain{intégrité} de Picard-Igusa, \add{et le th.\ de \emph{Néron-Severi} pour les surfaces proj.\ non sing., et \uncertain{résolution} des surfaces, de Abhyankar} \uncertain{suffit} à entraîner (vi) que si $X/S$ est propre et de prés.\ finie sur \struck{$k$} $S$ quasi-compact, les groupes de N.S. des fibres géométriques sont de \add{type fini} \struck{\ill{}} « \uncertain{uniforme} » \struck{\uncertain{et} \ill{}} (i.e.\ groupe de torsion d'ordre \uncertain{majoré}, et rang \uncertain{majoré}). [Dans le cas $X/S$ projectif, cela \uncertain{se} \uncertain{tire} des \uncertain{démarches} de Kleiman, sans utiliser aucun fourbis]. On récupère aussi ainsi (vii).

\uncertain{TSVP}

\page{71}
Ces résultats impliquent-ils le « critère de l'équivalence » sous la forme (ii) ? Qualitativement : « nettoyer » la situation, on \uncertain{est} ramené au cas où $\underline{\mathrm{Pic}}_{X/S}$, $\underline{\mathrm{Pic}}_{Y/S}$ sont représentables, $\underline{\mathrm{Pic}}^{\tau}_{X/S}$ et $\underline{\mathrm{Pic}}^{\tau}_{Y/S}$ aussi et plats et \uncertain{formés} \uncertain{d'ouverts}, et à examiner $\underline{\mathrm{NSL}}_{X/S} \to \underline{\mathrm{NSL}}_{Y/S}$. On \uncertain{tire} de \uncertain{ce} \uncertain{qui} précède, c'est que c'est un \emph{monomorphisme}. Il n'est \uncertain{pas} du tout clair pourquoi il serait de type fini. En fait, ma démonstration prouve que les morphismes
$\underline{\mathrm{Pic}}_{X/S} \to \underline{\mathrm{Pic}}_{Y/S}$, $\underline{\mathrm{NSL}}_{X/S} \to \underline{\mathrm{NSL}}_{Y/S}$, sont génériquement au dessus de $S$ des morphismes quasi-affines (peut-être \uncertain{même} affines). Moyennant (ii), on montre d'ailleurs que (ii) et (iii) sont essentiellement équivalents [le critère en termes du pol.\ de Hilbert, ou au choix, en termes de $D^2\cdot H^{n-2}$ et $D\cdot H^{n-1}$]. Je ne \uncertain{vois} pas que cela sorte du topos Kleiman.\note{la lecture « topos » est douteuse ; peut-être « type »}

\section{Lettre à Murre, 18 juillet 1962}
\note{Pp.~72--74 : copie dactylographiée au net, paginée par lui -2-, -3- ; pp.~75--78 : état antérieur de la même lettre, avec ses corrections à la main et à la machine, et signé. Les deux états sont donnés, puisque leurs différences sont l'information. La lettre est en anglais et reste en anglais.}

\page{72}
July 18, 1962

My dear Murre,

I recently had some thought on finiteness conditions for Picard preschemes, and substantially improved on the results stated in the last section of my last Bourbaki talk. The main result stated there for a simple projective morphism with connected geometric fibers (namely that the pieces $\underline{\mathrm{pic}}^{P}_{X/S}$ are of finite type over $S$) has been extended by Mumford to the case where instead of $f$ simple we assume only $f$ flat with integral geometric fibers, (at least if these are normal). Using his result (the proof of which is quite simple and beautiful), I could get rid of the normality assumption, and even (as in theorem 4.1. of my talk) restrict to the consideration of the two first non trivial coefficients of the Hilbert polynomials. The key results for the reduction are the following (the proofs being very technical, and rather different for (i) and (ii), except that (ii) uses (i) to reduce to the normal case; moreover (ii) uses Mumford's result and the equivalence criteria as developped in my last Seminar):

(i) Let $X$, $Y$ be proper over $S$ noetherian, let $f\colon X \to Y$ be a \emph{surjective} $S$-morphism, assume for simplicity of the statement that the Picard preschemes exist, the $f^{\cdot}\colon \underline{\mathrm{Pic}}_{Y/S} \to \underline{\mathrm{Pic}}_{X/S}$ is of finite type (and in fact affine if $S$ is the spectrum of a field), i.e.\ a subset $M$ of $\underline{\mathrm{Pic}}_{Y/S}$ is quasi-compact iff its image in $\underline{\mathrm{Pic}}_{X/S}$ is.

(ii) The same conclusion holds for a canonical immersion $X \to Y$, if $Y/S$ is projective with fibers all components of which are of dimension $\geqslant 3$, and if $X$ is the sub-scheme of zeros of a section over $Y$ of an invertible sheaf $\underline{L}$ ample with respect to $S$.

A connected result is that for any $X/S$ proper, and integer $n \neq 0$, the $n$.\,th power homomorphism in the Picard prescheme is of finite type.

I tell you about this, namely (i), because of the method of proof, involving of course considerations of non flat descent. The fact that I do not have any good effectivity criterion does not hamper, by just recalling what the effectivity of a given descent datum means. Now it turns out that by a slightly more careful analysis of the situation, one can prove the following theorem, of a type very close to the one you have proved recently, and to some you still want to prove as I understand it.

\page{73}
\emph{Theorem} Let $S$ be an integral noetherian scheme, $X$ and $X'$ proper over $S$ and $f\colon X' \;X$\note{la flèche manque dans la frappe} a surjective $S$-morphism, look at the corresponding homomorphisme for the Picard functors $f^{\cdot}\colon \underline{\underline{\mathrm{Pic}}}_{X/S} \to \underline{\underline{\mathrm{Pic}}}_{X'/S}$. Assume: a) the existence problems A and B defined below for $X/S$ has always a solution (this is certainly true when $X/S$ is projective). b) the morphism $f_s\colon X'_s \to X_s$ induced on the generic fiber is a morphism of descent, i.e.\ $\underline{O}_{X_s} \to f(\underline{O}_{X'_s}) \rightrightarrows h(\underline{O}_{X''_s})$ is exact. Then, provided we replace $S$ by a suitable non empty open set, the homomorphism $f^{\cdot}$ is representable by a quasi-affine morphism, more specifically in the factorisation of $f^{\cdot}$ via the functor representing suitable descent data, $f^{\cdot} = vu$ with $u$ affine and $v$ a monomorphism (as you well know), $v$ is in fact representable by \add{finite} direct sums of immersions.\note{« direct » est corrigé à la main sur « divest »}

\emph{Corollary} Without assuming b), but instead in a) allowing $X/S$ to be replaced by suitable other schemes $X_1$ finite over $X$, the same conclusion holds, namely $f^{\cdot}$ is representable by quasi-affine morphisms.

This follows from the theorem, using a suitable factorisation of $f$. For instance, using Chow's lemma and the Main existence theorem in my first talk on Picard schemes, one gets:

\emph{Corollary 2} Assume $X/S$ proper satisfies the condition a') for every $X'$ finite over $X$, there exists a non empty open subset $S_1$ of $S$ such that problem A for $X'$ $S_1$ has allways solution \struck{for problems A and B} (this condition is satisfied if $X/S$ is projective). Then provided we replace $S$ by a suitable $S_1$ non empty and open, $\underline{\mathrm{Pic}}_{X/S}$ exists, is separated, its connected components are of finite type over $S$.

\emph{N.B.} The proof does not give any evidence towards the fact that in the theorem, one could replace "quasi-affine" by "affine". This is true however over a field, bacause a quasi-affine algebraic group is affine! It would be interesting to have a couterexample, say, over a ring of dimension 1 such as $k[t]$\note{les crochets sont ajoutés à la main}, $X$ and $X'$ projective and simple over $S$ and $X' \to X$ birational, or alternatively, $X$ and $X'$ projective and normal over $S$, and $f\colon X' \to X$ finite. A counterexample in the latter case would of course provide a counterrexample to the effectivity problem for a finite morphism raised in my first talk on descent …
\note{Un trait vertical en marge longe ce N.B.}

"Problem A" is the following: given $X/S$ and Module $F$ on $X$, to represent the functor on the category of $S$-preschemes taking any $S'/S$ into a one-element or into the empty set, occording as to whether $F'$ on $X'/S'$ is flat with respect to $S'$ or not, where $X' = X \times_S S'$, $F' = F \times_S S'$.

\page{74}
Given $X/S$, we say that "Problem A for $X/S$ has always a solution" if for every coherent $F'$ on some $X'/S'$, the previous functor on \add{$(\mathrm{Sch})$}$/S'$ is representable by a $S'$-scheme of finite type. The main step in my proof of existence of Hilbert schemes shows that this condition is satisfied when $X/S$ is projective. In the proof, essential use is made of the Hilbert polynomial, in fact we get a solution as a disjoint sum of subschemes of $S$ corresponding to various Hilbert polynomials. Still I would expect that the functor is representable as soon as $X/S$ is proper. In view of the application we have in mind here, it would be sufficient (for any integral $S$) to find in $S$ a non empty open set $S_1$ such Problem A has always a solution for $X_1 = X \times_S S_1$ over $S_1$. To prove this weeker existence result, it is well possible that a reduction to the projective case is possible, using Chow's lemma and some induction on the relative dimension perhaps. I also would expect that a proof wil be easier when working over a complete noetherian local ring, hence the case of a general noetherian local ring by flat descent. And it is well possible that, putting together two such partial results, a proof of the existence in general could be obtained. (I met with such difficulties allready time ago in a very analoguous non projective existence problem, which beside I did not solve so far!). This problem A has been met also by Hartshorne (A Harvard Student), but I doubt he will work seriously on it. Thus I now wrote you in the hope you may be interested to have a try on this problem. As a general fact, our knowledge of non projective existence theorems is exceedingly poor, and I hope this will change eventually.

Sincerely yours.

A. GROTHENDIECK

\page{75}
\add{July 18, 1962}

My dear Murre,

I recently had some thought on finiteness conditions for Picard preschemes, and substantially improved on the results stated in the last section of my last Bourbaki talk. The main result stated there for a simple \struck{morphism (namely that} \add{projective morphism} with connected geometric fibers (namely that the pieces $\underline{\mathrm{Pic}}^{P}_{X/S}$ are of finite type over $S$) has been extended by Mumford to the case where \struck{the fibers} instead of $f$ simple we assume only $f$ flat with integral geometric fibers, \add{(at least if these are normal).} Using his result (the proof of which is quite simple and beautiful), I could get rid of the normality assumption, and even (as in theorem 4.1. of my talk) restrict to the consideration of the two first non trivial coefficients of the Hilbert polynomials. The key results for the reduction are the following (the proofs being very technical, and rather different for (i) and (ii), except that (ii) uses (i) \struck{and Mumford's} to reduce to the normal case; moreover (ii) uses Mumford's result and the equivalence criteria as developped in my last Seminar):

(i) Let $X$, $Y$ be proper over $S$ noetherian, let $f\colon X \to Y$ be a \add{surjective} $S$-morphism, \struck{\ill{}} assume for simplicity of the statement that the Picard preschemes exist, the $f^{\cdot}\colon \underline{\mathrm{Pic}}_{Y/S} \to \underline{\mathrm{Pic}}_{X/S}$ is of finite type (and in fact affine if $S$ is the spectrum of a field), i.e.\ a subset $M$ of $\underline{\mathrm{Pic}}_{Y/S}$ is quasi-compact iff its image in $\underline{\mathrm{Pic}}_{X/S}$ is.

(ii) The same conclusion holds for a canonical immersion $X \to Y$, if $Y/S$ is projective with fibers all components of which are of dimension $\geqslant 3$, and if $X$ is the sub-scheme of zeros of a section over $Y$ of an invertible sheaf $\underline{L}$ ample with respect to $S$.

A connected result is that for any $X/S$ proper, and integer $n \neq 0$,

\page{76}
the $n$.th power homomorphism in the Picard prescheme is of finite type.

I tell you about this, namely (i), because of the method of proof, involving of course considerations of non flat descent. The fact that I do not have any good effectivity criterion does not hamper, by just recalling what the effectivity of a given descent datum means. \struck{By} Now it turns out that by a slightly more careful analysis of the situation, one can prove the following theorem, of a type very close to the one you have proved recently, and \add{to} some you still want to prove as I understand it.

\emph{Theorem} Let $S$ be an integral noetherian scheme, $X$ and $X'$ proper over $S$ and $f\colon X' \to X$ a surjective $S$-morphism, look at the corresponding homomorphisme for the Picard functors $f^{\cdot}\colon \underline{\underline{\mathrm{Pic}}}_{X/S} \to \underline{\underline{\mathrm{Pic}}}_{X'/S}$. Assume: a) the existence problem\struck{s} A \struck{and B}\note{« and B » est rayé à la main, mais la copie au net p.~73 le garde} defined below for $X/S$ has always a solution (this is certainly true when $X/S$ is projective). b) the morphism $f_s\colon X'_s \to X_s$ induced on the generic fiber is a morphism of descent, i.e.\ $\underline{O}_{X_s} \to f_{*}(\underline{O}_{X'_s}) \rightrightarrows h_{*}(\underline{O}_{X''_s})$ is exact. Then, provided we replace $S$ by a suitable non empty open set, the homomorphism $f^{\cdot}$ is representable by \struck{a} quasi-affine morphism\struck{s}, more specifically \struck{that} in the factorisation of $f^{\cdot}$ via the functor representing suitable descent data, $f^{\cdot} = vu$ with $u$ affine \struck{\ill{}} and $v$ a monomorphism (as you well know), $v$ is in fact representable by \add{finite direct sums of} immersions.
\marginal{\uncertain{$S$} \ill{} \uncertain{needed} ?}\note{note au crayon en marge, face à « provided we replace $S$ … », que ce passage souligne}

\emph{Corollary} Without assuming b), but instead in a) allowing $X/S$ to be replaced by suitable other schemes $X_1$ finite over $X$, the same conclusion holds, namely $f^{\cdot}$ is representable by \struck{finite} quasi-affine morphisms.

This follows from the theorem, using a suitable factorisation of

\page{77}
$f$. For instance, using Chow's lemma and the Main existence theorem in my first talk on Picard schemes, one gets:

\emph{Corollary 2} Assume $X/S$ proper satisfies the condition \struck{a') for every \ill{}} a') for every $X'$ finite over $X$, there exist a non empty open subset $S_1$ of $S$ such that \add{problem A for} $X'|S_1$ has allways solution \struck{for problems A and B.} (this condition is satisfied if $X/S$ is projective). Then \struck{there exists} provided we replace $S$ by a suitable $S_1$ non empty and open, $\underline{\mathrm{Pic}}_{X/S}$ exists, is separated, its connected components are of finite type over $S$.

N.B. The proof does not give any evidence towards the fact that in the theorem, one could replace "quasi-affine" by "affine". This is true however over a field, bccause a quasi-affine algebraic group is affine! It would be interesting to have a couterexample, say, over a \struck{\ill{}} ring of dimension 1 such as $k[t]$, $X$ and $X'$ \add{projective and} simple over $S$ and $X' \to X$ birational, \struck{or} alternatively, \struck{\ill{}} $X$ and $X'$ projective and normal over $S$, and $f\colon X' \to X$ finite. A counterexample in the latter case would of course provide a couterexample to the effectivity problem raised in my first talk on descent for a finite morphism …
\marginal{(\,) impossible if $X$ smooth!}\note{à la main, en marge, face au contre-exemple ; « simple » est cerclé et relié à « projective and » ajouté au-dessus ; « for a finite morphism » est encerclé et ramené avant « raised »}

"Problem A" is the following: given $X/S$ and \struck{any} Module $F$ on $X$, to represent the functor on the category of $S$-preschemes taking any $S'/S$ into \struck{the empty} into a one-element \add{or into the empty} set, according as to wheter $F'$ on $X'/S'$ is flat with respect to $S'$ \add{or not}, where $X' = X \times_S S'$, $F' = F \otimes_S S'$. Given $X/S$, we say that "Problem A \add{for $X/S$} has always a solution" if for every coherent $F'$ on \struck{$X$} \add{some $X'/S'$}, the previous functor \add{on $(\mathrm{Sch})/S'$} is representable by a \add{$S'$-}scheme of finite type. The main step in my proof of existence of Hilbert schemes shows that this condition is satisfied when $X/S$ is projective.

\page{78}
In the proof, essential use is made of the Hilbert polynomial, in fact we get a solution as a disjoint sum of subschemes of $S$ corresponding to various Hilbert polynomials. Still I would expect that the \struck{\ill{}} functor is representable as soon as $X/S$ is proper. In view of the application we have in mind here, it would be sufficient (for any integral $S$) to find in $S$ a non empty \struck{\ill{}} open set $S_1$ such Problem A has always a solution for $X_1 = X \times_S S_1$ over $S_1$. To prove this weeker existence result, it is well possible that a reduction to the projective case is possible, using Chow's lemma and some induction on the relative dimension perhaps. I also would expect that a proof will be easier when working over a complete noetherian local ring, hence the case of a general \add{noetherian} \struck{local} ring by flat descent. And it is well possible that, putting together two such partial results, a proof of the existence in general could be obtained. (I met with such difficulties allready time ago in a very analoguous non projective existence problem, which beside I did not solve so far!). This problem A has been met also by Hartshorne (a Harvard student), but I doubt he will work seriously on it. Thus \add{I} now wrote you in the hope you may be interested to have a try on this problem. As a general fact, our knowledge of non projective existence theorems is exceedingly poor, and I hope this will change eventually.

Sincerely yours

\add{A Grothendieck}

\section{Extension du dual d'une VA et faisceaux inversibles à connexion}
\note{p.~79 : feuillet de couverture orange, de sa main : « Extension du \uncertain{dual} d'une VA et faisceaux inversibles à connexion. »}

\page{80}
\note{Cette page ouvre une suite qui se poursuit au-delà du lot.}
\begin{tikzcd}
X \arrow[d, "f"'] \\
S
\end{tikzcd}

\marginal{plat de prés.\ finie — en partie biffé}

\[
(*)\qquad \underline{\mathrm{Div}}_{X/S} \simeq \underline{R}^{*}_{X/S} / \underline{O}^{*}_{X} \xrightarrow{\ df/f\ } \underline{R}_{X/S}(\underline{\Omega}^{1}_{X/S}) / \underline{\Omega}^{1}_{X/S}
\]
\struck{\ill{}} d'où une extension [des couples $(D, \omega)$ tels que $d\log D \equiv \omega \ (\underline{\Omega}^{1}_{X/S})$]
\[
0 \to \underline{\Omega}^{1}_{X/S} \to \underline{E}_{X/S} \to \underline{\mathrm{Div}}_{X/S} \to 0
\]
\marginal{N.B. L'hom.\ $(*)$ est injectif si $S$ de car.\ nulle et $f$ normal \struck{\ill{}} (donc de prés.\ finie …)}
et par suite une suite exacte
\[
0 \to f_{*}(\underline{\Omega}^{1}_{X/S}) \to f_{*}(\underline{E}_{X/S}) \to f_{*}(\underline{\mathrm{Div}}_{X/S}) \xrightarrow{\ \partial\ } R^{1}f_{*}(\underline{\Omega}^{1}_{X/S})
\]
\note{sous $f_{*}(\underline{\Omega}^{1}_{X/S})$, un signe $=$ vertical renvoie à $\underline{\omega}_{X/S}$}
\marginal{à vérifier que c'est l'homomorphisme bien connu déduit de $\underline{\mathrm{Pic}}_{X/S} \to R^{1}f_{*}(\underline{\Omega}^{1}_{X/S})$, c'est trivial}
ou encore
\[
0 \to \underline{\omega}_{X/S} \to f_{*}(\underline{E}_{X/S}) \to f_{*}(\underline{\mathrm{Div}}_{X/S})^{w} \to 0
\]
où
\[
f_{*}(\underline{\mathrm{Div}}_{X/S})^{w} = \mathrm{Ker}\bigl(\partial\colon f_{*}(\underline{\mathrm{Div}}_{X/S}) \to R^{1}f_{*}(\underline{\Omega}^{1}_{X/S})\bigr)
\]

Considérons
\[
\underline{P}_{X/S} = f_{*}(\underline{R}^{*}_{X/S}) / f_{*}(\underline{O}^{*}_{X}) \hookrightarrow f_{*}(\underline{\mathrm{Div}}_{X/S})
\]
On a, du moins si $X \to S$ est \emph{quasi-projectif} \struck{\ill{}}
\[
f_{*}(\underline{\mathrm{Div}}_{X/S}) / \underline{P}_{X/S} \simeq \underline{\mathrm{Pic}}_{X/S}
\]
\struck{d'autre part, \ill{}} et \uncertain{l'on} \uncertain{voit}
\[
\begin{aligned}
f_{*}(\underline{\mathrm{Div}}_{X/S})^{w} / \underline{P}_{X/S} &= \underline{\mathrm{Pic}}^{w}_{X/S} \subset \underline{\mathrm{Pic}}_{X/S} \\
&= \mathrm{Ker}\bigl(\underline{\mathrm{Pic}}_{X/S} \to R^{1}f_{*}(\underline{\Omega}^{1}_{X/S})\bigr)
\end{aligned}
\]
* On définit la \uncertain{flèche} \uncertain{oblique}\note{fin de la page ; la phrase se poursuit au-delà du lot}

\end{document}
